%% kapa: exam=1 round=24 year=2013 period=1 form=A subject=civil_law subject_ko=민법 seq=6 number=6
%% source: Q-Net 감정평가사 기출문제 (cst00309) artlSeq=5073802 file=extracted/1008/5073802/2013년 제24회 감정평가사 1차시험 기출문제/1교시(A형).hwp
%% conversion: kordoc 4.24.51 -> tools/gen1.py (auto); review=proofread-sonnet; flags=none
\begin{question}{6}
물건에 관한 설명으로 옳지 않은 것은? (다툼이 있으면 판례에 의함)
\begin{choices}
  \item 1동의 건물의 일부도 독립성이 인정되면 하나의 부동산이 될 수 있다.
  \item 부동산의 일부를 전세권의 객체로 할 수 있다.
  \item 입목에 관한 법률에 의해 등기한 수목의 집단은 저당권의 객체가 될 수 없다.
  \item 독립된 부동산으로서의 건물이라고 하기 위하여는 최소한의 기둥과 지붕 그리고 주벽(主壁)을 갖추어야 한다.
  \item 타인의 토지 위에 권원없이 경작한 농작물이 성숙하였다면 그 소유권은 경작자에게 귀속된다.
\end{choices}

eviewflag{choices 3-4 overlap in rhwp render; only part verified}
\end{question}
\answer{3}
